\studySubsection{practice-calc-12}{练习库：高等数学第 12 讲 一元函数积分学的应用（三）——物理应用与经济应用}

\problemAnchor{MATH1-CALC-0059}
\begin{problemBox}
\begin{problemMeta}
\textbf{题目编号：} \problemRef{MATH1-CALC-0059}\quad
\textbf{答案：} \answerRef{MATH1-CALC-0059}\quad
\textbf{来源：} Codex 原创 / K110 深度讲义配套练习（非真题）
\end{problemMeta}
\begin{enumerate}
  \item[\textbf{A}] 物体沿直线从 \(x=0\) 移到 \(x=2\,\mathrm m\)，
  同向力为 \(F(x)=3x^2+2\,\mathrm N\)。求所做的功。
  \item[\textbf{B}] 半径为 \(1\,\mathrm m\)、高为 \(4\,\mathrm m\) 的
  竖直圆柱形水箱装满密度为 \(\rho\) 的液体。求把全部液体泵到水箱上沿
  所做的功，重力加速度记为 \(g\)。
  \item[\textbf{C}] 一块竖直等腰三角形薄板的顶点在水面，底边水平且位于
  深度 \(h\)，底边长 \(b\)。水密度为 \(\rho\)，求一侧总液体压力。
  \item[\textbf{M}] 均匀方向的细杆长 \(L\)，其线密度
  \[
  \lambda(x)=\lambda_0\left(1+\frac xL\right),\qquad0\le x\le L.
  \]
  质量为 \(m\) 的质点位于杆延长线上，距近端 \(a>0\)。求杆对质点的
  万有引力大小。
\end{enumerate}
\end{problemBox}

\problemAnchor{MATH1-CALC-0060}
\begin{problemBox}
\begin{problemMeta}
\textbf{题目编号：} \problemRef{MATH1-CALC-0060}\quad
\textbf{答案：} \answerRef{MATH1-CALC-0060}\quad
\textbf{来源：} Codex 原创 / K111 深度讲义配套练习（非真题）
\end{problemMeta}
\begin{enumerate}
  \item[\textbf{A}] 求 \(f(x)=\sin x\) 在 \([0,\pi]\) 上的平均值。
  \item[\textbf{B}] 细杆占据 \([0,2]\)，线密度
  \[
  \rho(x)=
  \begin{cases}
  1,&0\le x\le1,\\
  2,&1<x\le2.
  \end{cases}
  \]
  求总质量和质心。
  \item[\textbf{C}] 求均匀薄片
  \[
  0\le y\le1,\qquad y^2\le x\le y
  \]
  的质心。
  \item[\textbf{M}] 细杆占据 \([0,1]\)，线密度
  \(\rho_\lambda(x)=1+\lambda x\)，其中 \(\lambda\ge-1\)。
  求质心 \(\bar x(\lambda)\)，并证明它随 \(\lambda\) 单调增加。
\end{enumerate}
\end{problemBox}
